\section{Relative adjunction inputs}\label{sec:relative-inputs}

The index reductions use uniform torsion on rationally connected bases;
the effective application uses full field generation over a big base.
We record the three exact results proved in the companion
\emph{Relative denominators and effective systems for log Calabi--Yau
fibrations} \cite{OpenAIRelative}. Their proofs are independent of the
canonical-index theorem proved here and of abundance after nonvanishing.

A generalized pair $(Z,B_Z+M_Z)$ includes a rational b-divisor
$\mathbf M$ that is the pullback of a nef rational Cartier divisor
on a higher projective model. The divisor $M_Z$ is its trace on $Z$.
The expression ``$p\mathbf M$ b-Cartier'' means that this nef divisor
has a Cartier multiple $pM_W$ on one determination $W$, whose
pullbacks give the data on all higher models. Generalized discrepancies
are computed after subtracting the nef data on such a model. The
following result concerns an exact presentation by rational divisors.

\noindent\textit{Source: \cite[Theorem~\RDenominatorNumber]{OpenAIRelative}.}
\begin{theorem}[Uniform relative denominators]
\label{thm:relative-denominators}
Fix a finite set $\Phi\subset[0,1]\cap\Q$ and an integer $1\leq d\leq4$.
There are positive integers $p_0\mid p$, with $p_0$ clearing $\Phi$, and a
DCC set $\mathcal B\subset[0,1]\cap\Q$, depending only on $d$ and $\Phi$,
with the following property.

Let $(X,B)$ be a projective lc pair of dimension $d$, with effective rational
boundary having coefficients in $\Phi$.  Suppose that $f:X\to Z$ is a
contraction onto a projective base $Z$, with $\dim Z>0$, and
\[
 K_X+B\sim_\Q f^*D
\]
for a $\Q$-Cartier divisor $D$ on $Z$.  There exist
$\psi\in\C(X)^*$ and a $\Q$-Cartier divisor $D_Z\sim_\Q D$ such that
\begin{equation}\label{eq:relative-exact-presentation}
 K_X+B+\frac1{p_0}\Div(\psi)=f^*D_Z,
 \qquad D_Z=K_Z+B_Z+M_Z.
\end{equation}
Here $B_Z$ is effective with coefficients in $\mathcal B$,
$(Z,B_Z+M_Z)$ is generalized lc, and $\mathbf M$ is b-nef with
$p\mathbf M$ b-Cartier.  If $(X,B)$ is klt, the generalized pair on $Z$ is
generalized klt.
\end{theorem}

The constant $p_0$ trivializes the geometric generic adjoint, while
$p$ controls the moduli divisor on a determination. These are actual
divisor statements, not numerical equivalences. A varying coefficient
introduced only to construct a contraction will not enter the fixed
coefficient set when this theorem is applied with boundary zero.

\noindent\textit{Source: \cite[Proposition~\REffectiveNumber]{OpenAIRelative}.}
\begin{proposition}\label{prop:effective-base}
Under the hypotheses of Theorem~\ref{thm:relative-denominators}, suppose in
addition that $D$ is big.  There is a positive integer $m=m(d,\Phi)$ such
that, for every positive multiple $l$ of $m$, the complete divisorial
system
\[
 \left|\left\lfloor l(K_X+B)\right\rfloor\right|
\]
is nonempty and generates precisely the subfield $\C(Z)\subset\C(X)$.
Its rational map is therefore birationally equivalent to $f$ and to the
Iitaka fibration of $K_X+B$.
\end{proposition}

Here the complete floor system generates the entire relatively
algebraically closed base field, including its algebraic part. The
conclusion holds for every positive multiple of the stated degree.

\noindent\textit{Source: \cite[Proposition~\RCTorsionNumber]{OpenAIRelative}.}
\begin{proposition}\label{prop:rc-torsion}
Fix an integer $1\leq d\leq 4$, a DCC set
$I\subset [0,1]\cap\Q$, and a positive integer $p$.
There exists a positive integer $\ell=\ell(d,I,p)$ with the following
property.  Let $Z$ be a normal projective complex variety of dimension $d$
with a rationally connected smooth projective resolution.  Suppose that
$(Z,B+M_Z)$ is generalized klt, that $B\geq 0$ has coefficients in $I$,
and that the b-nef rational b-divisor $\mathbf M$ has $p\mathbf M$
b-Cartier.  If
\[
 D=K_Z+B+M_Z\sim_{\Q}0,
\]
then $\ell D$ is an integral principal divisor.  Consequently, for every
integer $a\geq1$, the divisor $a\ell D$ is principal and
\[
 h^0\bigl(Z,\OO_Z(a\ell D)\bigr)=1.
\]
The same integer can be chosen for all dimensions in any fixed subset of
$\{1,2,3,4\}$.
\end{proposition}

The last statement will be used directly with zero boundary and zero
nef data on a klt variety with rationally connected resolution. It also
applies to the generalized base of the first statement. After taking a
common multiple with $p_0$, its principal divisor pulls back to an integral
principal canonical multiple on the total space. This is the only
uniform-torsion step in Lemma~\ref{lem:index-moving-divisor}.
